
\documentclass{report}

\input{preamble}
\input{macros}
\input{letterfonts}

\title{\Huge{Algebra Lineare}\\Sistemi Informativi Aziendali}
\author{\huge{Caberlotto Giovanni}}
\date{}

\begin{document}

\maketitle
\newpage% or \cleardoublepage
% \pdfbookmark[<level>]{<title>}{<dest>}
\pdfbookmark[section]{\contentsname}{toc}
\tableofcontents
\pagebreak

\chapter{Vettori e Matrici} % (fold)
\label{chap:Vettori e Matrici}

\section{Concetto di vettore}
\subsection{Punti di vista differenti}

La pietra portante alla basa di tutta l'agebra lineare è il vettore, ci sono tre idee distinte ma correlate di cosa sia un vettore. dal punto di vista della fisica i vettori sono frecce nello spazio, e ciò che definisce un vettore è la lunghezza e la direzione in cui punta, ma fintanto che queste due quantità sono le stesse lo si può spostare e rimane lo stesso vettore. I vettori che si trovano nel piano sono bidimesionali, i vettori che invece si trovano nello spazio sono tridimensionali. Dal punto di vista informatico un vettore è una lista ordinata di elementi, per esempio mettiamo caso di star raccogliendo dati riguardo ai prezzi di alcune case, e le uniche quantità che ci interessano sono la metratura e il prezzo, possiamo moddelare quindi ogni singola casa con una coppia di numeri, il primo rappresenta la metratura e il secondo il prezzo l'ordine di questi due ha importanza per mettere a confronto ogni singola coppia, in gergo stiamo moddellando le case come vettori bidimensionali dove vettore non sta altro che a significare "lista", e ciò che lo rende bidimensionale è il fatto che la lunghezza di questa è 2. Dal punto di vista matematico invece si cerca di generalizzare entrambi i punti di vista appena elencati, affermando che un vettore può essere qualsiasi cosa per cui vi sia una sensata nozione di addizione o moltiplicazione per un numero i dettagli di questo punto di vista sono molto astratti

\subsection{Vettori nel piano bidimensionali}

Per comprendere i vettori nel piano, è essenziale avere familiarità con il concetto di sistema di riferimento cartesiano. In questo contesto, ogni punto del piano può essere associato univocamente a una coppia ordinata di numeri reali $(x, y)$, corrispondenti alle sue coordinate.

Un vettore nel piano è rappresentato come un segmento orientato che collega l'origine $(0, 0)$ al punto considerato. Questo segmento può essere visualizzato come una freccia con direzione e lunghezza. Le componenti del vettore sono i numeri reali $(x, y)$ associati al punto finale del segmento.

Se introduciamo un sistema di riferimento cartesiano il piano si può indentificare con l'insieme $\mathbb{R}^2$ delle coppie ordinate di numeri reali.


Ad ogni punto è possibile associare un vettore che possiamo pensare come un segmento orientato avente per estremi l'origine e il punto stesso

Ne segue quindi che possiamo "indentificare" il vettore $\vec{op}$ con la coppia ordinata $(x,y)$ delle coordinate del punto $p$ e scrivere $\vec{v} = (x,y)$ invece di $\vec{v} = \vec{op}$ dove $(x,y)$ sono detti componenti scalari e $|\vec{v}|$ è definito norma o modulo

\mprop{$|\vec{v}| =$ Lunghezza di $\vec{op} = \sqrt{x*2 + y^2}$}

\newpage

\subsubsection{Operazioni tra vettori}

\dfn{Somma vettoriale}{Dati due vettori $\vec{v} = (x,y)$ e $\vec{w} = (x,y)$ la loro somma è un vettore avente per componenti la somma delle componenti di $\vec{v}$ e $\vec{w}$}

\mprop{}{$\vec{v} + \vec{w} = (x_v + x_w,y_v + y_w)$}

\qs{}{Siano $\vec{v} = (1,3)$ e $\vec{w} = (4,1)$ due vettori nel piano. determinare $\vec{v} + \vec{w}$}

\sol{}{$\vec{v} + \vec{w} = (1+4,3+1) = (5,4)$}

Se volessimo dare un significato geometrico all'operazione appena svolta, vi sono più modi per rappresentarla:

% Aggiungere rappresentazione grafica 

\cor{}{Il modulo (norma) del vettore somme è minore uguale dei moduli dei 2 vettori sommati
\\
Disuguaglianza triangolare: $|\vec{v} + \vec{w}| \le |\vec{v}| + |\vec{w}|$
}

\dfn{Moltiplicazione per uno scalare}{Dati il vettore $\vec{v} = (x,y)$ ed il numero reale $\lambda$ il vettore $\lambda\vec{v}$ ha per componenti quelle del vettore $\vec{v}$ moltiplicate per $\lambda$}

\mprop{}{$\lambda\vec{v} = (\lambda x, \lambda y)$}

\qs{}{Sia $\vec{v} = (1,2)$, determinare $2\vec{v}$ e $-\vec{v}$}

 \sol{}{$2\vec{v} = (2*1,2*2)=(2,4)$
 \\ $-\vec{v} = (-1*1,-1*2) = (-1,-2)$}

 Se volessimo dare un significato geometrico all'operazione appena svolta:

% aggiungere rappresentazione geometrica

\cor{}{Nella moltiplicazione per $\lambda$ : il modulo (norma) viene moltiplicato per $|\lambda|$ la direzione resta la stessa, il verso resta uguale se $\lambda > 0$ mentre si inverte se $\lambda < 0$}

\subsubsection{Prodotto scalare e angolo tra vettori}

% chapter Vettori e Matrici (end)

\end{document}
